\begin{lemma}[Weak gradient of a positive part]
Let $n\ge1$, let $q\ge1$ and $c$ be real, let $\Omega\subseteq\R^n$ be open with
$|\Omega|<\infty$, and let $u\in W^{1,q}(\Omega)$ with weak gradient $g$
(\noderef{MemW1pOn}).  Then
\[
(u-c)^+=\max\{u-c,\,0\}\in W^{1,q}(\Omega)
\quad\text{with weak gradient}\quad
x\longmapsto g(x)\,\chi_{\{u>c\}}(x)
\]
(\noderef{MemW1pOn}), where $\chi_{\{u>c\}}$ is the indicator of
$\{x:u(x)>c\}$, so that the weak gradient is $g(x)$ when $u(x)>c$ and $0$
otherwise.
\end{lemma}
\begin{proof}
By \noderef{MemW1pOn} the assertion has two halves: the test-function identity
\[
\int_\Omega(u-c)^+\,D_i\varphi\,dx=-\int_\Omega g^i\,\chi_{\{u>c\}}\,\varphi\,dx
\tag{1}
\]
for every $\varphi\in C^\infty(\R^n)$ with compact support contained in $\Omega$
and every $i\in\{1,\dots,n\}$ (that is, \noderef{HasWeakGradientOn} for the pair
$((u-c)^+,\ g\chi_{\{u>c\}})$), and the integrability statements.

\emph{Step 0: integrability.}  Pointwise $0\le(u-c)^+\le|u|+|c|$, and
$u\in L^q(\Omega)$ by \noderef{MemW1pOn} while the constant $|c|$ lies in
$L^q(\Omega)$ because $|\Omega|<\infty$; hence $(u-c)^+\in L^q(\Omega)$.  Also
$|g\chi_{\{u>c\}}|\le|g|$ pointwise and $|g|\in L^q(\Omega)$, so
$|g\chi_{\{u>c\}}|\in L^q(\Omega)$.  Finally, as in the companion lemma,
$q\ge1$ and $|\Omega|<\infty$ give $g^i\in L^1(\Omega)$ for each $i$, which we
use for the dominated convergence below.

\emph{Step 1: a $C^1$ approximation of the positive part.}  For $\varepsilon>0$
define $F_\varepsilon\colon\R\to\R$ by
\[
F_\varepsilon(t)=
\begin{cases}
0, & t\le0,\\[2pt]
\dfrac{t^2}{2\varepsilon}, & 0\le t\le\varepsilon,\\[6pt]
t-\dfrac{\varepsilon}{2}, & t\ge\varepsilon .
\end{cases}
\]
The three formulas agree at the junctions ($0=0$ at $t=0$, and
$\varepsilon^2/(2\varepsilon)=\varepsilon/2=\varepsilon-\varepsilon/2$ at
$t=\varepsilon$), so $F_\varepsilon$ is well defined and continuous, and the
piecewise derivative is
\[
F_\varepsilon'(t)=
\begin{cases}
0, & t\le0,\\[2pt]
\dfrac{t}{\varepsilon}, & 0\le t\le\varepsilon,\\[6pt]
1, & t\ge\varepsilon ,
\end{cases}
\]
which again agrees at the junctions ($0=0$ at $t=0$ and $1=1$ at
$t=\varepsilon$).  Hence $F_\varepsilon$ is of class $C^1$ on $\R$, and
\[
0\le F_\varepsilon'(t)\le1\qquad\text{for every }t\in\R .
\tag{2}
\]
Two further properties, both immediate from the formulas:
\[
\bigl|F_\varepsilon(t)-t^+\bigr|\le\frac{\varepsilon}2
\qquad\text{for every }t\in\R ,
\tag{3}
\]
because the difference is $0$ for $t\le0$, equals
$t-\frac{t^2}{2\varepsilon}\in[0,\varepsilon/2]$ for $0\le t\le\varepsilon$ (it
is increasing there, with value $\varepsilon/2$ at $t=\varepsilon$), and equals
$\varepsilon/2$ for $t\ge\varepsilon$; and
\[
F_\varepsilon'(t)\longrightarrow\chi_{\{t>0\}}
\qquad\text{as }\varepsilon\downarrow0,\ \text{for every }t\in\R .
\tag{4}
\]
For (4): if $t\le0$ then $F_\varepsilon'(t)=0$ for every $\varepsilon>0$, and if
$t>0$ then $F_\varepsilon'(t)=1$ as soon as $\varepsilon<t$.  There is no
exceptional level set: the convergence in (4) holds at \emph{every} real $t$,
including $t=0$, which is the reason for using this piecewise $F_\varepsilon$
rather than a smoothing of $t\mapsto t^+$ whose derivative at $0$ would tend to
an unspecified value in $[0,1]$.

\emph{Step 2: applying the $C^1$ chain rule.}  Put
$H_\varepsilon(t)=F_\varepsilon(t-c)$.  Then $H_\varepsilon$ is $C^1$ with
$H_\varepsilon'(t)=F_\varepsilon'(t-c)$, so $|H_\varepsilon'|\le1$ everywhere by
(2).  By \noderef{WeakGradientC1Comp}, applied with $F=H_\varepsilon$, $L=1$ and
the same $n,q,\Omega,u,g$ --- whose hypotheses are exactly the ones we have,
namely $\Omega$ open with $|\Omega|<\infty$, $u\in W^{1,q}(\Omega)$ with weak
gradient $g$, and $H_\varepsilon$ of class $C^1$ with derivative bounded by $1$
---
\[
H_\varepsilon\circ u\in W^{1,q}(\Omega)
\quad\text{with weak gradient}\quad
H_\varepsilon'(u)\,g .
\]
In particular, by \noderef{HasWeakGradientOn}, for every admissible $\varphi$
and every $i$,
\[
\int_\Omega F_\varepsilon(u-c)\,D_i\varphi\,dx
=-\int_\Omega F_\varepsilon'(u-c)\,g^i\,\varphi\,dx .
\tag{5}
\]

\emph{Step 3: letting $\varepsilon\downarrow0$ in (5).}  Fix $\varphi$ and $i$,
and put $K=\operatorname{supp}\varphi$, a compact subset of $\Omega$; both
integrands in (5) vanish off $K$.

Left-hand side.  By (3), applied pointwise at $t=u(x)-c$,
\[
\Bigl|\int_\Omega\bigl(F_\varepsilon(u-c)-(u-c)^+\bigr)D_i\varphi\,dx\Bigr|
\le\frac{\varepsilon}2\int_K|D_i\varphi|\,dx\longrightarrow0
\qquad(\varepsilon\downarrow0),
\]
the integral $\int_K|D_i\varphi|\,dx$ being finite because $D_i\varphi$ is
continuous and $K$ compact.  Hence the left-hand side of (5) converges to
$\int_\Omega(u-c)^+D_i\varphi\,dx$.

Right-hand side.  By (4), applied pointwise at $t=u(x)-c$, we have
$F_\varepsilon'(u(x)-c)\to\chi_{\{u(x)-c>0\}}=\chi_{\{u>c\}}(x)$ for
\emph{every} $x$, with no exceptional set.  Moreover, by (2),
\[
\bigl|F_\varepsilon'(u-c)\,g^i\,\varphi\bigr|\le\|\varphi\|_\infty\,|g^i|\,\chi_K ,
\]
which is integrable because $g^i\in L^1(\Omega)$ by Step 0 and $K\subseteq\Omega$.
The dominated convergence theorem therefore gives
\[
\int_\Omega F_\varepsilon'(u-c)\,g^i\,\varphi\,dx
\longrightarrow\int_\Omega\chi_{\{u>c\}}\,g^i\,\varphi\,dx .
\]

Passing to the limit on both sides of (5) yields (1).  Together with Step 0 this
is precisely the assertion that $(u-c)^+\in W^{1,q}(\Omega)$ with weak gradient
$g\chi_{\{u>c\}}$ (\noderef{MemW1pOn}).
\end{proof}
